% Lean module: MonadFormal.Approximation.Events
\subsection{Event meaning of the error}

\begin{theorem}[Event characterization and separation]
\label{thm:variation-events}
For probability measures on any measurable space,
\[
 d_{\mathrm{TV}}(\mu,\nu)
   =\sup_{B\text{ measurable}}|\mu(B)-\nu(B)|,
 \qquad d_{\mathrm{TV}}(\mu,\nu)=0\iff\mu=\nu.
\]
\end{theorem}
\begin{proof}
Indicator functions are admitted tests, which proves that the event
supremum is no larger than $d_{\mathrm{TV}}$.
For the reverse inequality, fix $f\in\mathcal T(X)$. The layer-cake
formula gives
\[
 \int f\,d\mu=\int_{(0,1]}\mu\{x:f(x)\ge t\}\,dt,
 \qquad
 \int f\,d\nu=\int_{(0,1]}\nu\{x:f(x)\ge t\}\,dt.
\]
The tail probabilities are measurable, bounded functions of $t$.
Subtract these identities. The absolute value of each integrand
difference is bounded by the event supremum; the integration interval
has length one. Hence the test difference obeys the same bound.
Taking the supremum over $f$ proves equality. Zero distance then forces
agreement on every measurable event, which is equality of measures;
the converse is diagonal vanishing.
\leanproof{MonadFormal.Approximation.Test.indicator,MonadFormal.Approximation.eventTV,MonadFormal.Approximation.eventTV_le_tv,MonadFormal.Approximation.Test.tail,MonadFormal.Approximation.Test.integral_eq_integral_tail,MonadFormal.Approximation.tv_eq_eventTV,MonadFormal.Approximation.tv_eq_zero_iff}
\end{proof}
